% !PRINTED 506
\weberpage{506}

%<<CONT
represent, which is in fact a group, as is to be seen from $\pi_1^{-1}\varkappa\pi_1\pi_1^{-1}\varkappa_1\pi_1 = \pi_1^{-1}\varkappa\varkappa_1\pi_1$. From this it follows at the same time that $\pi_1^{-1}Q\pi_1$ is \gesperrt{isomorphic} with $Q$ (\S.\,149).

The groups to which the conjugate functions (10) belong, namely
\weq{12}{Q,\quad \pi_1^{-1}Q\pi_1,\quad \pi_2^{-1}Q\pi_2,\ldots \pi_{j-1}^{-1}Q\pi_{j-1}}
we call \gesperrt{conjugate divisors} of $P$, or briefly conjugate groups\footnote{The name ``equally entitled subgroup'' is also in use for this.}.

If $\pi$ is an arbitrary permutation from $P$, and hence of the form $\varkappa\pi_r$, then the group $\pi^{-1}Q\pi$ is always contained among the conjugate groups (12), namely $= \pi_r^{-1}\varkappa^{-1}Q\varkappa\pi_r = \pi_r^{-1}Q\pi_r$.

The derivation of $\pi^{-1}Q\pi$ from $Q$ is also called a \gesperrt{transformation of the group $Q$ by $\pi$}, and $\pi^{-1}Q\pi$ a \gesperrt{group transformed from $Q$}.

For the formation of the permutations of the conjugate groups $\pi^{-1}Q\pi$ the following simple rule serves:

6. \gesperrt{One obtains the permutation $\pi^{-1}\varkappa\pi$ by carrying out the permutation $\pi$ in the cycles of $\varkappa$.}

To prove the rule, let, decomposed into its cycles (\S.\,153),
\[
\varkappa = (\alpha,\beta,\gamma\ldots)(\alpha',\beta',\gamma'\ldots)\ldots
\]
and let
\[
\pi =
\begin{pmatrix}
\alpha, & \beta, & \gamma & \ldots \alpha', & \beta', & \gamma' & \ldots \\
\alpha_\pi, & \beta_\pi, & \gamma_\pi & \ldots \alpha'_\pi, & \beta'_\pi, & \gamma'_\pi & \ldots
\end{pmatrix}.
\]
By $\pi^{-1}$, $\alpha_\pi$ passes into $\alpha$; by $\varkappa$, $\alpha$ passes into $\beta$; and by $\pi$, $\beta$ is carried into $\beta_\pi$. By $\pi^{-1}\varkappa\pi$, therefore, $\alpha_\pi$ passes into $\beta_\pi$. Since the same consideration is applicable to $\beta_\pi, \gamma_\pi\ldots$ and so forth, it follows that:
\[
\pi^{-1}\varkappa\pi = (\alpha_\pi,\beta_\pi,\gamma_\pi\ldots)(\alpha'_\pi,\beta'_\pi,\gamma'_\pi\ldots),
\]
by which the rule 6. is proved.

Of the decomposition into subsidiary groups we can make yet another application, to the proof of the theorem:

7. \gesperrt{The degree of a transitive permutation group of $m$ figures is always a multiple of $m$.}
